\documentclass[10pt,a4paper]{amsart}
\usepackage[margin=19mm]{geometry}
\usepackage{amsmath,amssymb,mathtools}
\usepackage[scr=rsfs,cal=euler]{mathalfa}
\usepackage{xfrac}
\usepackage[unicode,hidelinks,bookmarks=false]{hyperref}
\newcommand{\ie}{i.e.\ }
\newcommand{\cf}{cf.\ }
\newcommand{\resp}{respectively\ }
\newcommand{\Subsection}{\subsection}
\providecommand{\nobreakdash}{\nobreak\hbox{-}}
\setlength{\emergencystretch}{2em}
\multlinegap0pt
\theoremstyle{plain}
\newtheorem{thm}{Theorem}
\newtheorem{lem}{Lemma}
\newtheorem{prop}{Proposition}
\newtheorem{cor}{Corollary}
\theoremstyle{definition}
\newtheorem{definition}{Definition}
\newtheorem{ex}{Example}
\theoremstyle{remark}
\newtheorem{rem}{Remark}
\newcommand{\dx}{\,\text{dx}}
\title[Learned finite element solution operators]{Learning Nonlinear Finite Element Solution Operators using Multilayer Perceptrons and Energy Minimization: the energy identity for the nonlinear problem}
\author{Mats G. Larson, Carl Lundholm, Anna Persson}
\date{Source: arXiv:2412.04596v3 (CC BY 4.0)}
\begin{document}
\maketitle
\noindent\emph{Source and licence.} Learning Nonlinear Finite Element Solution Operators using Multilayer Perceptrons and Energy Minimization. Version arXiv:2412.04596v3 (CC BY 4.0), distributed under the Creative Commons Attribution 4.0 International licence, \url{http://creativecommons.org/licenses/by/4.0/}, as recorded in the arXiv licence metadata for this exact version and shown on the abstract page \url{https://arxiv.org/abs/2412.04596v3}. Full licence notice: licenses/arxiv-2412.04596-CC-BY-4.0.txt.\par\smallskip

\noindent\textit{Editorial scope.} Counted unit: the article's Lemma with source label \texttt{lem:nonlinenergyident} (source0.tex lines 192--206, the part of the article's lemma environment that carries the statement), the identity that expresses the difference of the nonlinear energy at an admissible function and at the discrete minimiser as the sum of the energy norm of their difference and an explicit remainder, delivered together with the article's complete original proof of it (source0.tex lines 207--230, the article's \texttt{proof} environment of 24 lines, which the article nests inside the lemma environment, so that the lemma's closing line, together with the immediately following commentary sentence of the article (lines 231--233), is retained as context). No mathematics is written, repaired or re-derived: statement and proof are the article's own text line for line, in the article's own notation (the nonlinear coefficient, the energy functional, the weak form, the derivative of the coefficient and the remainder), and no retained line is substituted, re-flowed or omitted. No further part of the article is required as context and none is retained: the only references inside the retained lines point at the two labels that the retained statement itself defines, so no cross-reference is left unresolved and no citation occurs in any retained line; consequently the wrapper carries no bibliography and no bibliography entry. Nothing else of the article is retained: its error estimate, its examples, its learning sections and its figures are omitted. Editorial formatting only: an AMS \texttt{amsart} preamble that restates the environments and the single macro the retained lines use, namely the article's declarations of Theorem, Lemma, Proposition, Corollary, Definition, Example and Remark, and its macro that typesets the integration symbol; the wrapper does not load or redistribute the article's own style files. The article is typeset in its own \texttt{article}-class format with its own layout and notation style files; no class or style file is shipped with this package and the wrapper is an amsart document. Numbering differs from the article, because the article numbers its theorem-like environments inside each section and only one environment is retained: the counted lemma prints as Lemma 1. No picture environment and no graphical inclusion occurs in any retained line, and the package delivers no image asset. One bundled source file, \texttt{sources/169686/source0.tex}, accompanies the delivered item as the exact source of every retained span. Every retained span is recorded in \texttt{delivery-evidence.json} with method \texttt{exact\_tex}.
\begin{lem}[An identity for nonlinear energy]\label{lem:nonlinenergyident}
Let $E(v) = \frac{1}{2}(A(v) \nabla v,\nabla v) - (f, v)$ with $A(v) \geq \alpha > 0$ and $\int_{\Omega} (A(v) + \frac{1}{2}A'(v)v) |\nabla v|^2 \dx > 0$ for all $v \in V_h$. Also let
\begin{equation}
(A(u_h) \nabla u_h, \nabla v) + \frac{1}{2}(A'(u_h)v \nabla u_h, \nabla u_h) = (f, v) \quad \forall v \in V_h \label{nlei_vf}
\end{equation}
and $\| v \|_{A(w)}^2 := (A(w) \nabla v, \nabla v)$. Then
	\begin{align}
	\boxed{
		\frac{1}{2}\| u_h - v \|_{A(u_h)}^2 + R(v, u_h) = E(v) - E(u_{h}) \quad \forall v \in V_h \label{nlei_mainres}
		}
	\end{align}
	where the remainder $R$ is
	\begin{align}
	R(v, u_h) = \frac{1}{2} \int_{\Omega} \big(A(v) - A(u_h)\big) | \nabla v |^2 - A'(u_h) (v - u_h) | \nabla u_h |^2 \dx \label{nlei_remainder}
	\end{align}

	\begin{proof}
		The right-hand side is
		\begin{equation}
		\begin{aligned}
			& \, E(v) - E(u_{h}) \\
			= & \, \bigg( \frac{1}{2}\| v \|_{A(v)}^2 - (f,v) \bigg) - \bigg( \frac{1}{2}\| u_h \|_{A(u_h)}^2 - (f,u_h) \bigg) \\
			= & \, \frac{1}{2}\| v \|_{A(v)}^2 - \frac{1}{2}\| u_h \|_{A(u_h)}^2 + (f,\underbrace{u_h - v}_{\in V_h}) \\
			= & \, \frac{1}{2}\| v \|_{A(v)}^2 - \frac{1}{2}\| u_h \|_{A(u_h)}^2 \\
			& + (A(u_h) \nabla u_h, \nabla (u_h - v)) + \frac{1}{2}(A'(u_h) (u_h - v) \nabla u_h, \nabla u_h) \\
		\end{aligned}
		\end{equation}
		where we have used \eqref{nlei_vf} with test function $u_h - v$ in the last step. Splitting the first term in the last row and adding and subtracting $\frac{1}{2}\| v \|_{A(u_h)}^2$, we have
		\begin{equation}
		\begin{aligned}
			& \, E(v) - E(u_{h}) \\
			= & \, \frac{1}{2}\| v \|_{A(v)}^2 - \frac{1}{2}\| u_h \|_{A(u_h)}^2 + \| u_h \|_{A(u_h)}^2 - (A(u_h) \nabla u_h, \nabla v) \\
			& + \frac{1}{2}(A'(u_h) (u_h - v) \nabla u_h, \nabla u_h) \pm \frac{1}{2}\| v \|_{A(u_h)}^2 \\
			= & \, \frac{1}{2}\| v \|_{A(u_h)}^2 + \frac{1}{2}\| u_h \|_{A(u_h)}^2 - (A(u_h) \nabla u_h, \nabla v) \\
			& + \frac{1}{2}(A'(u_h) (u_h - v) \nabla u_h, \nabla u_h) + \frac{1}{2}\| v \|_{A(v)}^2 - \frac{1}{2}\| v \|_{A(u_h)}^2 \\
			= & \, \frac{1}{2}\| u_h - v \|_{A(u_h)}^2 + \frac{1}{2}\| v \|_{A(v) - A(u_h)}^2 - \frac{1}{2}(A'(u_h) (v - u_h) \nabla u_h, \nabla u_h)
		\end{aligned}
		\end{equation}
		which is the left-hand side of \eqref{nlei_mainres} with $R(v, u_h)$ on norm-product form.
	\end{proof}

\end{lem}

We note that Lemma~\ref{lem:nonlinenergyident} is indeed a generalization of the linear case with $A$ only depending on the coordinates of the solution domain, since then $A'(v) = 0$. The remainder $R$ in \eqref{nlei_mainres} also vanishes in that case, since $A(v) - A(u_h) = A - A = 0$ and $A'(u_h) = 0$.

\end{document}
